% Step 3: Hyperparameter table for paper appendix
% All values extracted from training scripts

\begin{table*}[t]
\centering
\small
\begin{tabular}{@{}lcccc@{}}
\toprule
\textbf{Hyperparameter} & \textbf{Script Reproduction} & \textbf{M2M-100 Hybrid} & \textbf{M2M-100 Conservative} & \textbf{mBART-50} \\
\midrule
Base model & M2M-100 418M & M2M-100 418M & M2M-100 418M & mBART-50 611M \\
Optimizer & Adam & AdamW & AdamW & AdamW \\
Learning rate & 3e-5 & 3e-5 & 1e-5 & 3e-5 \\
LR scheduler & None (fixed) & Cosine + warmup & Cosine + warmup & Cosine + warmup \\
Warmup steps & --- & 1,000 & 1,000 & 500 \\
\addlinespace
Batch size (per GPU) & 16 & 8 & 8 & 4 \\
GPUs & 1 & 1 & 1 & 3$\times$A6000 \\
Gradient accumulation & --- & 36 & 36 & 24 \\
Effective batch size & 16 & 288 & 288 & 288 \\
Max epochs & 20 & 20 & 20 & 20 \\
\addlinespace
Dropout & default & 0.2 & 0.2 & 0.2 \\
Attention dropout & default & 0.05 & 0.05 & 0.05 \\
Activation dropout & default & 0.05 & 0.05 & 0.05 \\
Weight decay & --- & 0.1 & 0.1 & 0.1 \\
Label smoothing & --- & 0.15 & 0.15 & 0.15 \\
\addlinespace
Max sequence length & 128 & 112 & 112 & 112 \\
Eval period (steps) & 1,000 & 500 & 500 & 500 \\
Early stopping patience & --- & 15 & 15 & 15 \\
Mixed precision & No & Yes & Yes & Yes \\
Data upsampling & --- & 5$\times$ & 5$\times$ & 5$\times$ \\
\addlinespace
Best checkpoint (steps) & 43,040 & 4,000 & 11,000 & 2,000 \\
Best validation loss & 1.60 & 3.31 & 3.34 & 1.27 \\
\bottomrule
\end{tabular}
\caption{Training hyperparameters for all model configurations. The original \texttt{train.py} uses minimal regularization (no scheduler, no label smoothing, no explicit dropout tuning). Our retrained models use consistent regularization (effective batch 288, dropout 0.2, label smoothing 0.15) with learning rate as the primary differentiator: 3e-5 (aggressive, from original paper) vs.\ 1e-5 (conservative). ``---'' indicates the parameter was not explicitly set or is not applicable.}
\label{tab:hyperparameters}
\end{table*}
